\section{Honesty-immunity: why a benchmark carries zero \Validated{} tags by design}
\label{sec:immunity}

The companion method paper \citep{baweja2026cti} carries exactly one externally
\Validated{} row---a real-data overbound-coverage anchor---and is scrupulous about
where that word may and may not appear. This paper carries \emph{zero}. That is not
a weaker result; it is a categorically different, and we argue more useful, kind of
object. This section makes the argument, states the firewall precisely, and
describes the one oracle we do run and why it is not validation.

\subsection{A benchmark makes no accuracy claim, so it cannot overclaim}
\label{sec:noaccuracy}

The four honesty tags of this program partition statements by what could make them
wrong. A \Cited{} statement is wrong if the citation is misread. A \Proven{}
statement is wrong if the algebra is wrong---and is machine-checked. A \Modelled{}
statement is a representative construction that claims no external fidelity. A
\Validated{} statement is an \emph{accuracy} claim---that some modelled quantity
matches an external truth---and it is the only kind that can be optimistic against
reality, which is exactly why it demands an external oracle and the most conservative
scope language.

A benchmark inhabits none of the accuracy-claiming category. It does not assert
that any clock is accurate, that any protection level is tight, that any monitor is
safe, or that any fault rate matches the world. It asserts only, on a fixed
stimulus and a fixed rule: \emph{this monitor's declared protection level did, or
did not, bound its true error at these epochs}, and it tallies the outcome. There
is no external truth the tally could contradict, because the tally is a statement
about the monitor's own declarations, not about reality. We call this
\emph{honesty-immunity}: because the benchmark makes no accuracy claim, it cannot
be wrong about accuracy, and so there is nothing to validate. The absence of a
\Validated{} tag is therefore not a gap in the evidence---it is the formal
signature of an object that structurally cannot overclaim.

This is what makes such an object a stable reference to cite. A
citing author who relies on a monitor must, in the end, trust that monitor's
accuracy claim; a citing author who relies on \emph{this benchmark} need trust only
that a fixed, open, versioned rule was applied to a fixed, open menu---something
they can re-run themselves in one command (Section~\ref{sec:availability}). The
yardstick's authority comes not from an accuracy claim we ask readers to believe,
but from the reproducibility of a rule that claims nothing.

\subsection{The honesty firewall}
\label{sec:firewall}

We state the firewall as a set of hard rules, each mirrored by a machine-checked
honesty gate over the benchmark source (\code{tests/tib\_honesty\_gates.rs}).

\begin{enumerate}
  \item \textbf{Zero \Validated{} tags.} No row of this paper or its benchmark is
  \Validated{}. The gate \code{no\_validated\_accuracy\_claim\_in\_benchmark}
  forbids the token ``validated'' from co-occurring with any accuracy word in the
  source. The word appears in this manuscript only to be defined and then shown
  never to be earned.
  \item \textbf{Method citations, not conformance.} The Stanford diagram
  \citep{tossaint2007} and the WAAS MOPS \citep{rtcado229} are \Cited{} as the
  scoring \emph{method}. The gate \code{no\_forbidden\_overclaim\_phrases} forbids
  ``certified'', ``formally verified'', ``DO-178'', ``ECSS conformance'', and any
  technology-readiness token; \code{required\_citations\_present} requires the
  Stanford, WAAS, and Mizrahi (RFC~7384) attributions to be present. This crate is
  not DO-229, DO-178, or ECSS certified, and claims no readiness level.
  \item \textbf{Undetectable means absorbed, never detected.} The gate
  \code{undetectable\_faults\_never\_described\_as\_detected} enforces
  Theorem~\ref{thm:invariant} at the level of prose as well as type: no source line
  naming an undetectable fault may assert detection.
  \item \textbf{No worst-case holdover.} The phrase ``worst-case holdover'' is
  forbidden; the benchmark scores absorption of an offset, not a deterministic
  bound on a coast.
  \item \textbf{No misattribution.} The benchmark source and this manuscript carry
  no author attribution other than the named author; a gate forbids stray
  attribution tokens in the source.
\end{enumerate}

The firewall is not decoration. It is the reason a reader can take the
honesty-immunity claim at face value: every clause is checked by a test that fails
the build if violated, including a fail-loud guard that treats a moved or renamed
benchmark source as a firewall violation rather than a silently skipped scan.

\subsection{The one oracle, and why it is not validation}
\label{sec:oracle}

We run exactly one cross-check, and we are careful to name it correctly. The
scorer's arithmetic---does the Rust implementation of \eqref{eq:classify} and
\eqref{eq:rates} count the four regions correctly?---is verified against an
\emph{independent} \code{numpy} re-implementation of the same four-region rule on a
fixed synthetic sample set (\code{scripts/gen\_tib\_scorer\_reference.py} writing
\code{tests/fixtures/tib/reference.json}; the Rust side asserts agreement in
\code{tests/tib\_scorer\_reference.rs}). This is an \code{InternalConsistency}
oracle in the program's verification matrix, and it is \Modelled{}, not
\Validated{}, for a precise reason: it cross-checks the \emph{counting}, not any
accuracy. The two implementations agree that a given $(e,\PL,\AL)$ triple falls in
a given region; neither asserts that the region assignment matches any real-world
truth, because there is no real-world truth about a synthetic triple to match. The
fixture's own \code{NOTICE.md} states this in the same words: it is \emph{not} an
external dataset and \emph{not} a validated accuracy oracle; the benchmark is
honesty-immune, so only the counting is checked. The reference tally is fixed and
small---$n=80$ epochs classified as $40$ \Nominal{}, $20$ \MI{}, $10$ \HMI{}, and
$10$ \Unavail{}, with $\mathrm{hmi\_rate}=0.125>\IR$ so $\mathrm{coverage\_ok}$ is
false---and the Rust scorer reproduces every one of those integers and every rate
to \num{e-12}. That agreement is real and useful; it is just not, and is never
labelled, validation.
